\begin{table}[ht]
\centering
\caption{Planner output on all 30 challenge prompts (descriptive only)}
\label{tab:challenge_planner_detail}
\begin{tabular}{lllllllll}
\toprule
prompt\_id & challenge\_type & recommended\_tool & task\_type & outcome\_type & likely\_fallback\_by\_content\_match & api\_call\_succeeded & returned\_models & reasoning \\
\midrule
chal\_out\_of\_scope\_mixed\_01 & out\_of\_scope\_mixed & summary\_stats & unknown & unknown & False & True & ['gpt-5.4-2026-03-05'] & The user is asking for a weather forecast rather than a statistical analysis, so no supported statistical tool is appropriate; default to unknown with the fallback tool. \\
chal\_out\_of\_scope\_mixed\_02 & out\_of\_scope\_mixed & summary\_stats & unknown & unknown & False & True & ['gpt-5.4-2026-03-05'] & The request is creative and does not specify an analytical objective, estimand, outcome, or model; default to unknown with the fallback tool. \\
chal\_out\_of\_scope\_mixed\_03 & out\_of\_scope\_mixed & summary\_stats & unknown & unknown & False & True & ['gpt-5.4-2026-03-05'] & The user requests language translation rather than statistical analysis, so no analysis tool is applicable; defaulting to the fallback tool. \\
chal\_out\_of\_scope\_mixed\_04 & out\_of\_scope\_mixed & linear\_regression & regression & unknown & False & True & ['gpt-5.4-2026-03-05'] & The actionable statistical request is to fit a regression of outcome on treatment; the email request is non-statistical and out of scope. With unspecified outcome scale, a basic regression routing choice is linear\_regression. \\
chal\_out\_of\_scope\_mixed\_05 & out\_of\_scope\_mixed & summary\_stats & unknown & unknown & False & True & ['gpt-5.4-2026-03-05'] & The request is unrelated to statistical analysis or causal inference, so no analysis tool is appropriate; default to the fallback tool. \\
chal\_out\_of\_scope\_mixed\_06 & out\_of\_scope\_mixed & summary\_stats & descriptive & unknown & False & True & ['gpt-5.4-2026-03-05'] & The user requests a dataset summary, which is a descriptive task best routed to summary statistics. \\
chal\_out\_of\_scope\_mixed\_07 & out\_of\_scope\_mixed & summary\_stats & unknown & unknown & False & True & ['gpt-5.4-2026-03-05'] & The user asks to review a Python function for bugs, which is outside the supported statistical analysis routing tasks. Defaulting to unknown with the fallback tool. \\
chal\_out\_of\_scope\_mixed\_08 & out\_of\_scope\_mixed & summary\_stats & unknown & unknown & False & True & ['gpt-5.4-2026-03-05'] & This is a factual lookup question rather than a statistical analysis request, so no specialized analysis tool is appropriate. \\
chal\_out\_of\_scope\_mixed\_09 & out\_of\_scope\_mixed & causal\_ate & causal\_effect & unknown & False & True & ['gpt-5.4-2026-03-05'] & The statistical request is a causal treatment effect question asking for the ATE of a drug on an outcome, so causal\_ate is the best match; the reminder request is non-analytic and not relevant for tool routing. \\
chal\_out\_of\_scope\_mixed\_10 & out\_of\_scope\_mixed & summary\_stats & unknown & unknown & False & True & ['gpt-5.4-2026-03-05'] & The request is not a statistical analysis task and does not specify data, outcomes, or a causal/regression/survival objective; default to unknown with a generic summary-oriented route. \\
chal\_underspecified\_01 & underspecified & causal\_ate & causal\_effect & unknown & False & True & ['gpt-5.4-2026-03-05'] & The request asks for the effect of a treatment on outcomes, which is a causal effect question best routed to an ATE estimation tool. \\
chal\_underspecified\_02 & underspecified & causal\_ate & causal\_effect & unknown & False & True & ['gpt-5.4-2026-03-05'] & The request is to determine whether treatment worked using treatment and outcome columns, which maps most directly to estimating a causal treatment effect. \\
chal\_underspecified\_03 & underspecified & causal\_ate & causal\_effect & unknown & False & True & ['gpt-5.4-2026-03-05'] & The user asks for the impact of a drug, which implies estimating a treatment effect comparing treated versus control; outcome type is not specified, so causal ATE is the best general routing choice. \\
chal\_underspecified\_04 & underspecified & survival\_adjusted\_curves & survival & time\_to\_event & False & True & ['gpt-5.4-2026-03-05'] & The request asks to compare two treatment groups over the study period, which implies follow-up over time and is best handled as a survival comparison. \\
chal\_underspecified\_05 & underspecified & causal\_ate & causal\_effect & unknown & False & True & ['gpt-5.4-2026-03-05'] & The request explicitly asks for causal analysis, so the appropriate routing is to estimate a treatment effect, most naturally the ATE, using a causal effect tool. \\
chal\_underspecified\_06 & underspecified & causal\_ate & causal\_effect & unknown & False & True & ['gpt-5.4-2026-03-05'] & The question asks whether drug X makes a difference versus not taking it, which is a treatment effect question best framed as estimating a causal average treatment effect rather than simple description. \\
chal\_underspecified\_07 & underspecified & causal\_ate & causal\_effect & unknown & False & True & ['gpt-5.4-2026-03-05'] & The request explicitly asks to adjust for confounding and interpret treatment group results, which is a causal treatment effect question best handled by an ATE estimation tool. \\
chal\_underspecified\_08 & underspecified & causal\_ate & causal\_effect & unknown & False & True & ['gpt-5.4-2026-03-05'] & The request asks for the difference between treated and untreated groups, which is best framed as an average treatment effect comparison rather than pure description. \\
chal\_underspecified\_09 & underspecified & causal\_ate & causal\_effect & unknown & False & True & ['gpt-5.4-2026-03-05'] & The request explicitly asks for a causal effect estimate with proper adjustment, which maps directly to estimating an ATE using a causal inference tool. \\
chal\_underspecified\_10 & underspecified & causal\_ate & causal\_effect & unknown & False & True & ['gpt-5.4-2026-03-05'] & The user asks whether the intervention helped and explicitly requests adjustment for baseline group differences, which implies a causal treatment effect estimation problem rather than simple descriptive comparison. \\
chal\_unsupported\_method\_01 & unsupported\_method & causal\_ate & causal\_effect & continuous & False & True & ['gpt-5.4-2026-03-05'] & The request is for a causal treatment effect with endogeneity addressed via an instrumental variable; the estimand is an IV-based causal effect of statin uptake on a continuous outcome. \\
chal\_unsupported\_method\_02 & unsupported\_method & causal\_ate & causal\_effect & unknown & False & True & ['gpt-5.4-2026-03-05'] & The request is for a causal effect at a policy eligibility threshold, which is a regression discontinuity estimand. Among available tools, causal\_ate is the closest match for estimating a treatment effect, though the design specifically relies on the age cutoff. \\
chal\_unsupported\_method\_03 & unsupported\_method & causal\_ate & causal\_effect & continuous & False & True & ['gpt-5.4-2026-03-05'] & The request is to estimate a causal policy effect using treated/control and pre/post structure, which corresponds to a difference-in-differences estimand and is best routed as a causal effect analysis. \\
chal\_unsupported\_method\_04 & unsupported\_method & causal\_ate & causal\_effect & continuous & False & True & ['gpt-5.4-2026-03-05'] & The request is to estimate a policy's causal effect using a synthetic control design on panel data, so the statistical goal is causal effect estimation with a continuous outcome. \\
chal\_unsupported\_method\_05 & unsupported\_method & causal\_ate & causal\_effect & unknown & False & True & ['gpt-5.4-2026-03-05'] & The request is for causal decomposition of an intervention effect into mediated and direct components, which is a causal effect task centered on treatment and outcome with adherence as mediator. Among available tools, causal\_ate is the closest match for causal effect estimation. \\
chal\_unsupported\_method\_06 & unsupported\_method & causal\_ate & causal\_effect & unknown & False & True & ['gpt-5.4-2026-03-05'] & The request is to estimate heterogeneous causal effects, which is a causal inference task focused on subgroup-specific treatment effects. A causal effect tool is the closest match because it uses treatment, outcome, and covariates to estimate treatment effects, even though the specific method requested is a causal forest. \\
chal\_unsupported\_method\_07 & unsupported\_method & causal\_ate & causal\_effect & time\_to\_event & False & True & ['gpt-5.4-2026-03-05'] & The request is a causal target trial emulation with time-varying treatment and mortality outcome. The estimand is a causal effect under dynamic treatment strategies; among available tools, causal\_ate is the closest match for causal effect estimation, though the full parametric g-formula requires longitudinal structure and time-varying confounders. \\
chal\_unsupported\_method\_08 & unsupported\_method & causal\_ate & causal\_effect & unknown & False & True & ['gpt-5.4-2026-03-05'] & The request is explicitly causal and asks for front-door identification of the exposure effect using a mediator under unmeasured exposure-outcome confounding, so the best routing is a causal effect estimation tool. \\
chal\_unsupported\_method\_09 & unsupported\_method & causal\_ate & causal\_effect & unknown & False & True & ['gpt-5.4-2026-03-05'] & The request is for a causal treatment effect estimated via propensity score nearest-neighbor caliper matching rather than weighting, so the appropriate routing is a causal effect tool that can estimate an effect after matching. \\
chal\_unsupported\_method\_10 & unsupported\_method & causal\_ate & causal\_effect & unknown & False & True & ['gpt-5.4-2026-03-05'] & The request is explicitly for a causal effect analysis using a marginal structural model with inverse probability of treatment and censoring weights for longitudinal time-varying treatment data, so the closest routing choice is a causal ATE tool. \\
\bottomrule
\end{tabular}
\end{table}
